% GENERATED by techreport/figs/emit_icl_category_tex.py -- do not edit by hand.
% Paired comparison: the three conditions score the same flaws.

\begin{table}[t]
\centering\small
\setlength{\tabcolsep}{5pt}\renewcommand{\arraystretch}{1.15}
\begin{tabular}{@{}p{3.9cm} c c c c c c c@{}}
\toprule
 & & \multicolumn{3}{c}{\textbf{Flaws found, with}} & & & \\
\cmidrule(lr){3-5}
\textbf{Flaw category} & \textbf{Flaws} & \textbf{no lessons} & \textbf{8} & \textbf{33} & \textbf{Change} & \textbf{95\% interval} & \textbf{Lessons} \\
\midrule
\textbf{order and timing} & 15 & 6 & 7 & 10 & +4 (+27) & $[+1, +27]$ & 3 \\
how many & 15 & 8 & 9 & 10 & +2 (+13) & $[-11, +24]$ & 2 \\
on-screen text & 55 & 36 & 46 & 43 & +7 (+13) & $[-0, +20]$ & 1 \\
\textbf{action or event occurs} & 177 & 124 & 139 & 143 & +19 (+11) & $[+5, +14]$ & 13 \\
camera and style & 13 & 11 & 12 & 12 & +1 (+8) & $[-5, +8]$ & 0 \\
attribute & 38 & 24 & 24 & 26 & +2 (+5) & $[-4, +10]$ & 2 \\
object identity & 42 & 33 & 36 & 34 & +1 (+2) & $[-1, +2]$ & 3 \\
physical motion & 12 & 8 & 7 & 8 & +0 (+0) & $[-14, +14]$ & 0 \\
spatial relation & 40 & 25 & 24 & 23 & -2 (-5) & $[-12, +6]$ & 1 \\
\bottomrule
\end{tabular}
\caption{Effect of the distilled lessons on each flaw category. A flaw is \emph{found} when the critic's findings cover it wholly or in part, as in \cref{tab:headline}. The three middle columns count that category's flaws found under no lessons, under $8$ lessons distilled from $25$ experience clips, and under $33$ from $117$; the change is given in flaws with percentage points in brackets, since several categories carry few enough flaws that a rate alone would overstate the movement. Because the three conditions score the same flaws the comparison is paired, and the interval on the change is computed from the flaws that switched status. \textbf{Bold} marks the two categories whose interval excludes zero. The last column counts the lesson families added between $25$ and $117$ experience clips whose name refers to that category.}
\label{tab:iclcategory}
\end{table}
